\section{Introduction}
\label{sec:01_Introduction}

Aerodynamic design balances a trade-off between the demand for physical fidelity and the need for rapid design iteration. High-fidelity Computational Fluid Dynamics (CFD) simulations, such as those based on the Reynolds-Averaged Navier-Stokes (RANS) equations, offer unparalleled accuracy in predicting complex flow phenomena. However, their significant computational cost makes them impractical for design-space exploration or real-time analysis. Conversely, data-driven surrogate models, especially deep learning architectures like Convolutional Neural Networks (CNNs), can provide near-instantaneous predictions but are entirely dependent on the quality and availability of the high-fidelity data used for their training and the robustness of their design.\\

This project develops a pipeline that integrates both methods. Inspired by the framework of \citet{Cuozzo2026AIDriven}, we first build a high-performance CFD solver to generate reference-quality aerodynamic data. We then leverage this data not only to train a custom CNN surrogate model but also to perform a systematic investigation into its architecture and training protocol. By establishing this workflow for clean symmetric airfoils, we made a base idea for future extensions to more complex iced geometries.\\

The document is structured to follow this approach, beginning with the development of the high-fidelity simulation framework. Section \ref{sec:02_Navier-Stokes_Simulation} (\textbf{Navier-Stokes Simulation}) details the extension of the \textbf{\lifex} \cpp finite element library to handle external aerodynamic flows. We describe the integration of the one-equation Spalart-Allmaras RANS model, the adaptation of the solver to two-dimensional domains, the implementation of specialized boundary conditions and the development of utilities for computing aerodynamic forces. This section also provides a comprehensive overview of the High-Performance Computing (HPC) workflow established on the CINECA Leonardo cluster.\\

Then, Section \ref{sec:03_Convolutional_Neural_Network} (\textbf{Convolutional Neural Network}), presents the design, implementation and training of our surrogate model. A key contribution of this work is the development of a CNN entirely from scratch in \cpp, designed to predict the lift coefficient from an airfoil's geometric representation along with the angle of attack and Reynolds number. We elaborate on the mathematical foundations of the network architecture, the physics-informed loss function that regularizes training and the implementation of distributed-memory parallelism using MPI, enabling efficient execution on HPC hardware.\\

Building upon this foundation, Section \ref{sec:04_Tuning} (\textbf{Tuning}) transitions from model implementation to rigorous optimization. This chapter details the comprehensive hyperparameter tuning campaign undertaken to identify the most accurate and stable model configuration. Using a five-fold cross-validation methodology, we systematically investigate the impact of network topology, optimizer selection, regularization techniques (including dropout and the physics-informed loss term), activation functions and learning rate schedules. This empirical analysis provides the evidence-based justification for the final production model's design.\\

Finally, Section \ref{sec:05_Conclusions} (\textbf{Conclusions}) synthesizes the findings from all stages of the project, evaluates the performance of the final optimized model and discusses the implications and future directions of this integrated framework.\\